% Section 6: ethics and the limits of this project.
% The limitation facts follow a check of this checkout:
%   autoencoders/ holds no training code and no weights;
%   rectified_flow/final_report/ and rectified_flow/out/ are untracked run artifacts,
%   so the raw run records, the metric CSV and the 96-frame sets are not in the repository;
%   no FOMM checkpoint is present (fomm/vox-cpk.pth.tar is ignored and absent);
%   the mean absolute pixel difference helper works on 8-bit pixel values.

\section{Ethics and Limitations}

\subsection{Ethics}
Face reenactment carries a clear risk of misuse. The same pipeline that makes a
light-hearted clip can put a real face on words the person never said, and a
viewer cannot separate the two cases by eye. The risk grows as soon as the
material leaves the team.

We answer that risk with provenance, with watermarks and with detection
research. The technical measures limit what a stranger can do with a published
file.

C2PA is a signed provenance and edit-history manifest. It records which tool
made a file, what changed it and when, and a signature lets a reader check that
record. It is not a stamp that exists only at capture time, because an editor
can add a further manifest entry for each change, and the history then travels
with the file. C2PA is also not a deepfake detector by itself. The manifest
reports where a file came from and how it changed; it does not judge whether the
content is real.

Watermarking is a separate measure. It hides a signal inside the pixels or the
audio, and the signal travels with the content. A watermark and a C2PA manifest
do different jobs, so a project can use both.

Research also reports detection cues. Eye blinking follows a rhythm that the
early face generators did not reproduce (Li et al., 2018). Remote
photoplethysmography reads a heartbeat signal from small colour changes in the
skin (Qi et al., 2020). Frequency analysis finds the traces that neural
upsampling leaves in the image spectrum (Frank et al., 2020). We treat these
cues as research examples only. This project ran no detection experiment and
trained no detector, so we report no measurement of our own here.

\subsection{Limitations}
Our results rest on artifacts that are not all in the repository. The
autoencoder training code and the trained weights are absent, so that training
cannot be repeated from this checkout. The Rectified Flow raw run records, the
metric CSV and the 96-frame sets are absent as well, so the numbers in
Section 4 cannot be re-derived here. The FOMM checkpoint is absent too, so the
training history of those weights cannot be verified; the released
configuration trains no adversarial loss, as Section 3 states.

Two explanations stay unverified. We ran no analysis of the dataset content, so
the braces explanation in Section 3 stays a hypothesis and not a finding. The
temporal metric thresholds are not calibrated, so a face change value has no
absolute reading; only the control level gives it a meaning.

One helper limits what one number can show. The script that reports the mean
absolute difference converts both images to 8-bit pixel values before the
subtraction, an 8-bit round trip. The value therefore cannot show a difference
finer than one 8-bit level.

These limits do not change any measured value. They set the size of each claim
we make, and we report only what the artifacts support.
